% ECONOMICS_PROBLEM: {"id": "I11-577", "title": "Steering patients toward higher-value care", "jel_code": "I11", "source_id": "OP-0444"}
% !TeX program = xelatex
\documentclass[11pt,a4paper]{article}
\usepackage[margin=23mm]{geometry}
\usepackage{fontspec}
\setmainfont{Latin Modern Roman}
\usepackage{amsmath,amssymb,mathtools}
\usepackage{xcolor,enumitem,booktabs,longtable,array,xurl,hyperref}
\definecolor{muted}{HTML}{53636C}
\hypersetup{colorlinks=true,urlcolor=blue,linkcolor=blue}
\setlength{\parindent}{0pt}
\setlength{\parskip}{5pt}
\newcommand{\R}{\mathbb R}
\newcommand{\E}{\mathbb E}
\newcommand{\N}{\mathbb N}
\newcommand{\ind}{\mathbf 1}
\newcommand{\Var}{\operatorname{Var}}
\newcommand{\Cov}{\operatorname{Cov}}
\newcommand{\supp}{\operatorname{supp}}
\DeclareMathOperator*{\argmin}{arg\,min}
\DeclareMathOperator*{\argmax}{arg\,max}
\newcommand{\entrymeta}[3]{{\sffamily\small\color{muted}\textbf{\texttt{#1}}\quad|\quad #2\par #3\par}}
\newcommand{\Problemref}[1]{\texttt{#1}}
\newcommand{\JELref}[2]{\texttt{#2} (JEL \texttt{#1})}
\begin{document}
\section*{Steering patients toward higher-value care}
\textbf{Problem ID:} \texttt{I11-577}\quad \textbf{JEL:} \texttt{I11}\par
% BEGIN ECONOMICS STATEMENT
\label{problem:OP-0444}
{\sffamily\small\color{muted}Original subject group: Health economics\par}
\label{definitions:problem:30007726}
\textbf{Audit classification:} Background; proposed extension. Checked source identity and appropriate agenda/background support; expanded intervention, units, population, definedness and causal restrictions. Exact targets and variants are editorial.
\subsection*{Definitions and precise research target}
\textit{Editorial formalization.} Consider privately insured adults referred for elective joint replacement in one regional health-care market. Let \(z\in\{0,1\}\) indicate an offer combining provider quality information, scheduling assistance, and a stated travel subsidy. Let \(J_i(z)\) be the chosen provider, \(H_i(z)\) one-year health improvement, \(C_i(z)\) real treatment cost, and \(T_i(z)\) patient time and travel costs. For monetary health valuation \(\lambda>0\), define
\[\tau_W=E[\lambda\{H_i(1)-H_i(0)\}-\{C_i(1)-C_i(0)\}-\{T_i(1)-T_i(0)\}]-k,\]
where \(k\) is per-eligible-person program resource cost. Provider quality is estimated independently of treatment assignment, with risk adjustment specified before the intervention. Identify \(\tau_W\) through randomized offers and report uptake and provider-choice effects without excluding patients who decline assistance. At larger market saturation, allow queues, prices, and provider capacity to respond; a small trial need not identify these responses. Determine which component of the offer matters only when factorial variation or credible mechanism restrictions permit.

Let \(H_i(z)\) be incremental quality-adjusted life-years over one year relative to the predeclared common clinical baseline, \(C_i(z)\) real treatment/aftercare resource cost, and \(T_i(z)\) monetized patient time and travel burden. Travel-subsidy payments are transfers; the travel resource cost and subsidy administration cost enter once. \(k\) is expected incremental program resource cost per eligible patient and cannot overlap \(C_i\) or \(T_i\). Baseline eligibility precedes assignment, and nonreceipt/operation cancellation remains part of the assigned regime. Finite outcome means, consistency, offer randomization and positivity identify the average intention-to-treat contrast under no spillovers. Provider-specific causal health effects require separate clinical risk/selection assumptions; independent estimation prevents mechanical reuse of treatment outcomes but does not establish validity of a quality measure. All expectations use a stated baseline population and probability law. Identification means every admissible data-generating model with the same observed-data law yields the same displayed target; otherwise report the identified set. Exact equations, horizons and assumptions are editorial, rather than source-given canonical conjectures.
\subsection*{Variants and assumption changes}
\begin{enumerate}
\item \textbf{Editorial component variant.} Factorially randomize information, scheduling and travel assistance, producing eight arms. Define mean net-health-benefit interactions.
\item \textbf{Editorial congestion variant.} Randomize market saturation and define \(H_i(z,s),\allowbreak C_i(z,s),\allowbreak T_i(z,s)\). Identify how queues and capacity change the isolated-offer contrast.
\end{enumerate}
\subsection*{References and source audit}
\textbf{1.} Inertia/selection policy questions verified; exact steering experiment and payment estimands are editorial. Locator: Draft 24 March 2014; Section 7, printed p.40; Sections 8–9. Full text or primary publication page inspected. URL: \url{https://kateho.scholar.princeton.edu/sites/g/files/toruqf4136/files/kateho/files/healthcare_io_3_24_14.pdf}.
\begin{enumerate}\raggedright
\item\hypertarget{definitions:source:30007726:1}{} Martin Gaynor; Kate Ho; Robert J. Town. 2014. \emph{The Industrial Organization of Health Care Markets}. Draft 24 March 2014; Section 7, printed p.40; Sections 8–9.
\url{https://kateho.scholar.princeton.edu/sites/g/files/toruqf4136/files/kateho/files/healthcare_io_3_24_14.pdf}.
DOI: \url{https://doi.org/10.3386/w19800}.
\end{enumerate}

\subsection*{Common conventions: Source mathematical conventions}

These conventions apply to each entry unless explicitly replaced.

\textbf{Models and identification.} A primitive model \(m\) specifies all probability laws, feasible choices, information, equilibrium and measurement rules used by its target. The admissible class \(\mathcal M\) is fixed independently of outcomes. If \(P_O(m)\) is its observable probability law and \(\tau(m)\) the target, its identified set at observable law \(P\) is
\[
 \mathcal I_\tau(P)=\{\tau(m):m\in\mathcal M,\ P_O(m)=P\}.
\]
Point identification means that this set is a singleton. An empty set signals incompatibility of data and assumptions. Coordinate bounds need not describe the joint set. Bounds are sharp only if every retained target is attainable or arbitrarily approachable by compatible models. Finite bounds require explicit restrictions on unobserved quantities; unrestricted confounding, hidden wealth, extrapolation or arbitrary equilibrium selection can yield uninformative sets.

\textbf{Causal interventions.} A potential outcome \(Y_i(p)\) is the outcome under a complete policy regime \(p\), including timing, financing, information and market spillovers. The target population and units are fixed before treatment. With interference, use the complete assignment vector or a justified exposure mapping; individual no-interference notation alone cannot identify an economy-wide intervention. A difference of mean potential outcomes does not require the cross-world joint distribution. Conditioning on post-treatment migration, survival, enrollment or employment changes the estimand and needs separate justification.

\textbf{Statistical statements.} An estimator, confidence set, forecast or test specifies a sampling law, model class and loss/coverage criterion. Uniform \((1-\alpha)\)-coverage means \(\inf_{m\in\mathcal M}\Pr_m\{\tau(m)\in C_n\}\geq1-\alpha\), or an explicitly asymptotic version. The whole parameter space trivially has coverage, so informativeness requires a stated width, risk or power objective. A forecasting improvement is relative to a named benchmark, information set and evaluation population.

\textbf{Equilibrium and optimization.} An equilibrium comprises feasible plans, optimality, beliefs where needed, budget constraints and all market/resource clearing. The primitives or a stated benchmark define each correspondence \(\mathcal E(p;m)\). Nonemptiness is assumed or proved, never inferred just from compact policy sets. A selected-equilibrium target uses a declared selection rule; a robust target quantifies over all permitted equilibria. Use suprema or infima unless attainment is established. An \(\varepsilon\)-optimal policy is feasible and within \(\varepsilon\) of the finite optimum value. Report an empty feasible set separately.

\textbf{Welfare and accounting.} Normative utility, interpersonal normalization, social weights, numeraire, discounting and the use of public receipts are primitives. Behavioral utility need not coincide with the welfare criterion. Household welfare includes ownership income and rents once; adding profits again to compensating variation generally double-counts them. A monetary equivalent is defined by an explicit transfer and price convention, with monotonicity and range conditions. Average effects, mechanism contributions, transfers and welfare losses are different targets. Infinite sums require convergence and dynamic feasible plans require terminal or no-Ponzi conditions.

\textbf{Source versions.} A working-paper date, revision date and journal publication year are distinguished. A DOI for a journal version is not automatically the DOI of an earlier manuscript. Locators refer to the stated version and printed pages where specified. A metadata-only check does not verify a claim about a full-text conclusion. Every entry's source subsection narrows the provenance claim accordingly.
% END ECONOMICS STATEMENT
\end{document}
